% Figure: a spark-ignition engine's brake torque and brake specific fuel
% consumption (BSFC) against engine speed at wide-open throttle. Torque peaks
% in the mid-range; BSFC bottoms (best efficiency) near the torque peak and
% rises at both low and high speed. The two curves share the speed axis; BSFC
% is on the right axis. Numbers are representative of a 1.6 L naturally
% aspirated engine, not a specific product.
\begin{figure}[htbp]
\centering
\begin{tikzpicture}
\begin{axis}[
  width=12cm, height=8cm,
  xlabel={engine speed (rpm)},
  ylabel={brake torque (N$\cdot$m)},
  xmin=800, xmax=6500, ymin=80, ymax=160,
  axis lines=left,
  tick label style={font=\scriptsize}, label style={font=\small},
  clip=false,
  axis y line*=left,
]
% Brake torque: rises to a mid-range peak near 3500 rpm, then falls.
\addplot[engblue, very thick, smooth] coordinates {
  (1000,108) (1500,120) (2000,131) (2500,140) (3000,146)
  (3500,148) (4000,146) (4500,141) (5000,133) (5500,122) (6000,108)
};
\node[font=\scriptsize, engblue, anchor=south] at (axis cs:3500,149)
  {brake torque (left axis)};
\addplot[engblack, only marks, mark=*, mark size=2pt] coordinates {(3500,148)};
\node[font=\scriptsize, anchor=north] at (axis cs:3500,146)
  {torque peak};
\end{axis}
% === second axis for BSFC, overlaid ===
\begin{axis}[
  width=12cm, height=8cm,
  xmin=800, xmax=6500, ymin=220, ymax=340,
  axis y line*=right,
  axis x line=none,
  ylabel={BSFC (g/kW$\cdot$h)},
  ylabel style={font=\small}, tick label style={font=\scriptsize},
]
% BSFC: bottoms near the torque peak (best efficiency), rises at both ends.
\addplot[engred, very thick, dashed, smooth] coordinates {
  (1000,298) (1500,278) (2000,262) (2500,251) (3000,245)
  (3500,243) (4000,246) (4500,253) (5000,264) (5500,282) (6000,308)
};
\node[font=\scriptsize, engred, anchor=north east] at (axis cs:6000,306)
  {BSFC (right axis)};
\addplot[engblack, only marks, mark=*, mark size=2pt] coordinates {(3500,243)};
\node[font=\scriptsize, anchor=south, engred] at (axis cs:3500,246)
  {best BSFC};
\end{axis}
\end{tikzpicture}
\caption[Engine torque and brake specific fuel consumption]{Brake torque
(blue, left axis) and brake specific fuel consumption (red dashed, right axis)
for a naturally aspirated spark-ignition engine at wide-open throttle. The
torque peaks in the mid-speed range where the breathing of the engine is best
matched to the valve and manifold tuning. The BSFC, the fuel mass burned per
unit of brake work, bottoms near the same speed: that minimum is the engine's
best operating point, where brake thermal efficiency is highest, because BSFC
and brake thermal efficiency are reciprocals once scaled by the fuel heating
value. A low BSFC of about $243\,$g/kW$\cdot$h corresponds to a brake thermal
efficiency near $34\,\%$ for a fuel of $43.5\,$MJ/kg lower heating value. The
rise in BSFC at low speed (friction and heat loss dominate the small work per
cycle) and at high speed (pumping and friction grow faster than work) is the
reason a transmission keeps the engine near the BSFC island during steady
cruise. Numbers are representative, not from a specific product.}
\label{fig:vol04ch03:engine-power-bsfc}
\end{figure}
